\titleframe

% Contoh Kode Slide
%\section{Introduction}
%\begin{frame}{\insertsectionhead}
%	\framesubtitle{A short introduction to Trigon}
%	\themename is a modern, elegant and versatile theme for Beamer, inspired by
%	the
%	\href{https://github.com/matze/mtheme}{\textsc{metropolis} theme} from Matthias
%	Vogelgesang.
%	\vfill
%	\themename comes with lots of nice extra features
%	\begin{itemize}
%		\item Multiple style variations for title, section and normal slides
%		\item Simple customization of theme colors
%		\item Lots of convenient options to tweak the design
%	\end{itemize}
%\end{frame}

% Contoh Kode Tabel
%  \begin{table}[H]
% \centering
% \caption{A nice table example}
% \begin{tabular}{@{} lccc @{}}
%    \toprule
%    & \textbf{Velocity} & \textbf{Angle}  & \textbf{Vertical force} \\
%    & $U$ & $\alpha$  & $F_z$ \\
%    & [m/s] & [$^\circ$]  & [N] \\
%    \midrule
%    2D simulation  & 9 & 2 & 9.23 \\
%    3D simulation  & 10.0 & 3 & 15.039 \\
%    Experiment A   & 11.31 & 2.5 & 13.2 \\
%    Experiment B   & 11.26 & 2.7 & 12.6 \\
%    Experiment C   & 11.33 & 2.47 & 13.6 \\
%    \bottomrule
%  \end{tabular}
%\end{table}

% Contoh Kode Gambar
%  \begin{figure}[ht!]
%	\begin{subfigure}[b]{0.3\textwidth}
%		\frame{\includegraphics[width=\textwidth]{layout_example-03.jpg}}
%		\caption*{plain}
%	\end{subfigure}
%	\hspace{\fill}
%	\begin{subfigure}[b]{0.3\textwidth}
%		\frame{\includegraphics[width=\textwidth]{layout_example-02.jpg}}
%		\caption*{style1}
%	\end{subfigure}
%	\hspace{\fill}
%	\begin{subfigure}[b]{0.3\textwidth}
%		\frame{\includegraphics[width=\textwidth]{layout_example-01.jpg}}
%		\caption*{style2 (default)}
%	\end{subfigure}
%\end{figure}

% Contoh Kode Blocks
% \begin{block}{Regular block}
%	Just a regular block
%\end{block}
%\begin{alertblock}{Alert block}
%	Some important thing
%\end{alertblock}
%\begin{exampleblock}{Example block}
%	No difference with regular block to avoid excessive distraction
%\end{exampleblock}

\setbeamertemplate{frame footer}{Alauddin Maulana Hirzan, M. Kom - Universitas Semarang}

\section{Tantangan dan Isu \textit{Internet of Things}}
\begin{frame}{\insertsectionhead}
	\framesubtitle{Teknologi adalah Teknologi}
	\justifying
	Teknologi apapun pasti memiliki yang namanya baik, buruk, kesempatan, dan ancaman (\textbf{\textit{Strength, Weakness, Opportunity, and Threat}}). Sehingga dalam implementasinya harus mempertimbangkan situasi dan kondisinya.
	\vfill
	\begin{block}{Contoh}
		Penggunaan monitoring di pegunungan memiliki tantangan dalam penggunaan baterai dan sinyal.
	\end{block}
\end{frame}

\begin{frame}{\insertsectionhead}
	\framesubtitle{Tantangan dan Isu dunia IoT}
	\justifying
	Berdasarkan \textit{IEEE Newsline - Three Major Challenges Facing IoT} : Teknologi, Bisnis, dan Warga
	\vfill
	  \begin{figure}[ht!]
			\begin{subfigure}[b]{1\textwidth}
					\frame{\includegraphics[width=\textwidth]{tantangan.png}}
			\end{subfigure}
		\end{figure}
\end{frame}

\begin{frame}{\insertsectionhead}
	\framesubtitle{Tantangan dan Isu dunia IoT - Keamanan \#1}
	\justifying
	IoT telah berubah menjadi masalah keamanan serius yang telah menarik perhatian perusahaan teknologi terkemuka dan lembaga pemerintah di seluruh dunia. Peretasan monitor bayi, lemari es pintar, termostat, pompa infus obat, kamera, dan bahkan radio di mobil Anda menandakan mimpi buruk keamanan yang disebabkan oleh masa depan IoT.
	\vfill
	\begin{block}{Info}
		Karena perangkat IoT terhubung dengan Internet maka akses dari pihak luar menjadi mudah.
	\end{block}
\end{frame}

\begin{frame}{\insertsectionhead}
	\framesubtitle{Tantangan dan Isu dunia IoT - Keamanan \#2}
	\justifying
	Masih tidak yakin?
	\vfill
	\begin{figure}[ht!]
		\begin{subfigure}[b]{1\textwidth}
			\frame{\includegraphics[width=\textwidth]{hacked.png}}
		\end{subfigure}
	\end{figure}
	\vfill
	Dari kulkas menjadi \textbf{Spy Robot}
\end{frame}

\begin{frame}{\insertsectionhead}
	\framesubtitle{Tantangan dan Isu dunia IoT - Keamanan \#3}
	\justifying
	Para hacker dapat mengubah benda-benda \textit{smart} yang memiliki teknologi komputasi menjadi sesuatu hal yang berbahaya bagi penggunanya. Sehingga dapat mencuri hal-hal tertentu seperti data-data sensitif atau digunakan untuk \textit{mining}
	\vfill
	\begin{figure}[ht!]
		\begin{subfigure}[b]{0.35\textwidth}
			\frame{\includegraphics[width=\textwidth]{bitcoin-mining.jpg}}
		\end{subfigure}
	\end{figure}
\end{frame}

\begin{frame}{\insertsectionhead}
	\framesubtitle{Tantangan dan Isu dunia IoT - Keamanan \#4}
	\justifying
	Dampak apa saja yang bisa muncul:
	\begin{itemize}
		\item Kehilangan Informasi Sensitif
		\item Gangguan Bisnis
		\item Memanipulasi Informasi di Perjalanan
	\end{itemize}
	\vfill
	\begin{block}{Info}
		Sebagian besar pengguna Internet di dunia menggunakan password yang sederhana dan mudah ditebak. Menggunakan password yang rumit membantu pengamanan IoT
	\end{block}
\end{frame}

\begin{frame}{\insertsectionhead}
	\framesubtitle{Tantangan dan Isu dunia IoT - Keamanan \#5}
	\justifying
	Lalu, \textbf{bagaimana mengamankan perangkat Internet of Things?} Ada beberapa cara yang bisa dilakukan untuk mengamankan perangkat IoT seperti:
	\vfill
	\begin{itemize}
		\item Memasang \textbf{Firewall}
		\item Memasang \textbf{Intrusion Detection System}
		\item Teknologi NAT untuk memisahkan jaringan luar dan dalam
		\item Komunikasi terenkripsi
		\item DLL
	\end{itemize}
\end{frame}

\begin{frame}{\insertsectionhead}
	\framesubtitle{Tantangan dan Isu dunia IoT - Konektivitas \#1}
	\justifying
	Model saat ini cukup untuk ekosistem IoT di mana puluhan, ratusan, atau bahkan ribuan perangkat terlibat. Tetapi ketika jaringan tumbuh untuk menggabungkan miliaran dan ratusan miliar perangkat, sistem terpusat akan berubah menjadi hambatan
	\vfill
	\begin{figure}[ht!]
		\begin{subfigure}[b]{0.4\textwidth}
			\frame{\includegraphics[width=\textwidth]{peta-internet.jpg}}
		\end{subfigure}
	\end{figure}
\end{frame}

\begin{frame}{\insertsectionhead}
	\framesubtitle{Tantangan dan Isu dunia IoT - Konektivitas \#2}
	\justifying
	\begin{figure}[ht!]
		\begin{subfigure}[b]{0.65\textwidth}
			\frame{\includegraphics[width=\textwidth]{survey.png}}
		\end{subfigure}
	\end{figure}
\end{frame}

\begin{frame}{\insertsectionhead}
	\framesubtitle{Tantangan dan Isu dunia IoT - Konektivitas \#3}
	\justifying
	Menurut Statista, ada sekitar 21,5 miliar perangkat yang saling terhubung di dunia. Jumlah mereka akan meledak di tahun-tahun berikutnya karena konsumsi internet meningkat dan gadget serta mesin baru memasuki pasar. 
	\vfill
	\begin{block}{Info}
		Kenaikan ini menjadi awal mula dari \textit{Internet of Everything} yang dimana perangkat-perangkat akan terhubung ke Internet.
	\end{block}
\end{frame}

\begin{frame}{\insertsectionhead}
	\framesubtitle{Tantangan dan Isu dunia IoT - Kompatibilitas dan Keawetan \#1}
	\justifying
	Kompatibiltias dan keawetan merupakan tantangan yang harus diikuti oleh teknologi IoT. Semakin banyak perangkat yang terhubung dalam jaringan IoT, maka kompatibiltias harus dipertimbangkan agar membuat perangkat tersebut tetap bisa dipakai sampai kapapnpun.
	\vfill
	\begin{block}{Info}
		\textbf{Firmware} atau Sistem Operasi yang mempengaruhi keawetan dan kompatibilitas sebuah perangkat. Perangkat lama belum tentu dapat dipasang Firmware/OS terbaru.
	\end{block}
\end{frame}

\begin{frame}{\insertsectionhead}
	\framesubtitle{Tantangan dan Isu dunia IoT - Kompatibilitas dan Keawetan \#1}
	\justifying
	Raspberry Pi OS membuat sebuah sistem operasi yang kompatible dengan perangkat-perangkat SOC mereka. Mulai dari versi \textbf{Zero} hingga Raspberry Pi 4, semua didukung oleh satu sistem operasi yang sama.
	\vfill
	\begin{figure}[ht!]
	\begin{subfigure}[b]{0.43\textwidth}
		\frame{\includegraphics[width=\textwidth]{rpi.jpg}}
	\end{subfigure}
\end{figure}
\end{frame}

\begin{frame}{\insertsectionhead}
	\framesubtitle{Tantangan dan Isu dunia IoT - Kompatibilitas dan Keawetan \#2}
	\justifying
	Beberapa dari teknologi ini pada akhirnya akan menjadi usang dalam beberapa tahun ke depan, secara efektif membuat perangkat yang mengimplementasikannya menjadi tidak berguna. Ini sangat penting, karena berbeda dengan perangkat komputasi generik yang memiliki masa pakai beberapa tahun, peralatan IoT (seperti lemari es pintar atau TV) cenderung tetap beroperasi lebih lama, dan harus dapat berfungsi bahkan jika produsennya keluar dari layanan.
\end{frame}

\begin{frame}{\insertsectionhead}
	\framesubtitle{Tantangan dan Isu dunia IoT - Kompatibilitas dan Keawetan \#3}
	\justifying
	Namun hal ini masih membahas di satu perangkat yang sama pabrikannya. Yang menjadi pertanyaan adalah Bagaimana kompatibilitas antar perangkat \textit{heterogeneous} dapat terjaga?
	\vfill
	Di dunia memiliki berbagai macam jenis perangkat seperti:
	\begin{multicols}{2}
		\begin{itemize}
			\item Arduino
			\item Orange Pi
			\item Banana Pi
			\item Rock Pi
			\item NVidia Jetson
			\item Raspberry Pi
		\end{itemize}
	\end{multicols}
\end{frame}

\begin{frame}{\insertsectionhead}
	\framesubtitle{Tantangan dan Isu dunia IoT - Standar \#1}
	\justifying
	Standar teknologi yang mencakup protokol jaringan, protokol komunikasi, dan standar agregasi data, adalah jumlah dari semua aktivitas penanganan, pemrosesan, dan penyimpanan data yang dikumpulkan dari sensor. Tantangan yang dihadapi adopsi standar dalam IoT:
	\vfill
	\begin{itemize}
		\item  Standar untuk menangani data tidak terstruktur: Data terstruktur disimpan dalam database relasional dan ditanyakan melalui SQL misalnya.
		\item  Keterampilan teknis untuk memanfaatkan alat agregasi yang lebih baru
	\end{itemize}
\end{frame}

\begin{frame}{\insertsectionhead}
	\framesubtitle{Tantangan dan Isu dunia IoT - Standar \#2}
	\justifying
	Karena sebagian besar perangkat IoT menggunakan WiFi, maka ada beberapa standar yang dipakai:
	\vfill
	\begin{multicols}{2}
		\begin{itemize}
			\item IEEE 802.11 (defunct)
			\item IEEE 802.11a
			\item IEEE 802.11b
			\item IEEE 802.11g
			\item IEEE 802.11n
			\item IEEE 802.11ac
			\item IEEE 802.11ax
		\end{itemize}
	\end{multicols}
	\vfill
\end{frame}

\begin{frame}{\insertsectionhead}
	\framesubtitle{Tantangan dan Isu dunia IoT - Standar \#3}
	\justifying
	Apa yang terjadi jika memakai jaringan dengan standar berbeda:
	\vfill
	\begin{quote}
		Misalnya, jika si A membawa pulang router 802.11ac baru. Router tersebut tidak akan mendukung langsung perangkat lama yang ada. Si A hanya bisa menerima beberapa manfaat router, seperti peningkatan jangkauan, tetapi koneksi perangkatnya hanya secepat standar Wi-Fi perangkat.
	\end{quote}
\end{frame}

\begin{frame}{\insertsectionhead}
	\framesubtitle{Tantangan dan Isu dunia IoT - Analisis dan Aksi Cerdas \#1}
	\justifying
	Tahap terakhir dalam implementasi IoT adalah mengekstraksi wawasan dari data untuk analisis, di mana analisis didorong oleh teknologi kognitif dan model yang menyertainya yang memfasilitasi penggunaan teknologi kognitif.
	\vfill
	\begin{figure}[ht!]
		\begin{subfigure}[b]{0.43\textwidth}
			\frame{\includegraphics[width=\textwidth]{smart-iot.png}}
		\end{subfigure}
	\end{figure}
\end{frame}

\begin{frame}{\insertsectionhead}
	\framesubtitle{Tantangan dan Isu dunia IoT - Analisis dan Aksi Cerdas \#2}
	\justifying
	Sebagian besar perangkat IoT nantinya akan terdiri dari triliunan perangkat yang berbeda-beda. Ada yang menggunakan mikrokontroler dan ada yang menggunakan \textit{System-on-Chip}. Sehingga dalam kekuatan komputasinya akan berbeda-beda. Perangkat cerdas memerlukan teknologi komputasi agar bisa memroses data tersebut menjadi informasi.
	\vfill
	\begin{figure}[ht!]
		\begin{subfigure}[b]{0.8\textwidth}
			\frame{
\includegraphics[width=\textwidth]{tflite.png}}
		\end{subfigure}
	\end{figure}
\end{frame}

\begin{frame}{\insertsectionhead}
	\framesubtitle{Tantangan dan Isu dunia IoT - Analisis dan Aksi Cerdas \#3}
	\justifying
	Karena itulah seluruh perangkat mikrokontroler tidak bisa diimplementasikan dengan teknologi AI, dan hanya bekerja secara murni sebagai penerima input atau sensor.
	\vfill
	\begin{figure}[ht!]
		\begin{subfigure}[b]{0.6\textwidth}
			\frame{\includegraphics[width=\textwidth]{sensor.png}}
		\end{subfigure}
	\end{figure}
\end{frame}

\begin{frame}{\insertsectionhead}
	\framesubtitle{Tantangan dan Isu dunia IoT - Analisis dan Aksi Cerdas \#4}
	\justifying
	Tantangan yang dihadapi adopsi tindakan cerdas dalam IoT
	\vfill
	\begin{itemize}
		\item Tindakan mesin dalam situasi yang tidak terduga
		\item Keamanan dan privasi informasi
		\item Interoperabilitas mesin
		\item Perilaku manusia yang mengubah perilaku
		\item Adopsi teknologi baru yang lambat
	\end{itemize}
\end{frame}

\begin{frame}{\insertsectionhead}
	\framesubtitle{Tantangan dan Isu dunia IoT - Bisnis}
	\justifying
	Garis besarnya adalah motivasi besar untuk memulai, berinvestasi, dan menjalankan bisnis apa pun. Tanpa model bisnis yang sehat dan solid untuk IoT akan menjadi masalah lain. Model ini harus memenuhi semua persyaratan untuk semua jenis e-commerce; pasar vertikal, pasar horizontal, dan pasar konsumen. Tetapi kategori ini selalu menjadi korban pengawasan peraturan dan hukum.
	\vfill
	\begin{block}{Info}
		Menjalankan bisnis berbasiskan IoT tidaklah mudah, antara regulasi pemerintah dan kebutuhan teknologi yang sangat ketat menghambat perkembangan IoT di negeri sendiri.
	\end{block}
\end{frame}

\begin{frame}{\insertsectionhead}
	\framesubtitle{Tantangan dan Isu dunia IoT - Warga \#1}
	\justifying
	Memahami IoT dari calon pelanggan dan regulator bukanlah tugas yang mudah karena alasan berikut:
	\vfill
	\begin{itemize}
		\item Tuntutan dan persyaratan pelanggan berubah terus-menerus.
		\item Pertumbuhan dan perkembangan teknologi perangkat
		\item Menemukan dan mengintegrasikan kembali fitur dan kemampuan yang harus dimiliki itu mahal dan membutuhkan waktu dan sumber daya.
		\item Penggunaan teknologi Internet of Things berkembang dan berubah
	\end{itemize}
\end{frame}

\begin{frame}{\insertsectionhead}
	\framesubtitle{Tantangan dan Isu dunia IoT - Warga \#1}
	\justifying
	\begin{itemize}
		\item Keyakinan Konsumen
		\item Kurangnya pemahaman atau edukasi oleh konsumen tentang praktik terbaik untuk keamanan perangkat IoT untuk membantu meningkatkan privasi.
	\end{itemize}
	\vfill
	IoT menciptakan tantangan unik untuk privasi, banyak yang melampaui masalah privasi data yang ada saat ini. Sebagian besar dari ini berasal dari mengintegrasikan perangkat ke dalam lingkungan kita tanpa kita secara sadar menggunakannya.
\end{frame}

\setbeamertemplate{background} 
{
	\includegraphics[width=\paperwidth,height=\paperheight]{thank-you.jpg}
}
\begin{frame}[plain]
\end{frame}